\chapter{Introducere}\label{ch:intro}
\pagestyle{fancy}

\section{Contextul lucrării}
\label{sec:intro-context}

Protezele robotice reprezintă unul dintre cele mai active domenii de cercetare la
intersecția dintre ingineria calculatoarelor, biomedicală și robotică. Capacitatea de
a reproduce mișcările naturale ale mâinii umane are implicații profunde atât în
domeniul medical, restabilirea funcționalității membrelor superioare pentru
persoanele cu dizabilități, cât și în cel industrial, unde mâinile robotice sunt
utilizate în manipularea obiectelor, asamblare de precizie și interacțiune
om-robot.

Un aspect esențial al acestor sisteme este \textit{controlul în timp real}. Mâna
umană posedă 27 DOF (\textit{Degrees of Freedom}, numărul de mișcări independente
posibile) și poate efectua mișcări cu o latență de ordinul milisecundelor. Pentru
a reproduce un comportament natural, un sistem robotic trebuie să proceseze datele
senzoriale, să determine pozițiile dorite ale articulațiilor și să genereze
semnalele de control în cadrul unor constrângeri temporale stricte. Orice latență
sau fluctuație în semnalele de comandă se traduce în mișcări sacadate, nepotrivite
cu gestul intenționat.

În paralel, domeniul recunoașterii gesturilor a cunoscut progrese semnificative
datorită dezvoltării rețelelor neuronale și a bibliotecilor de viziune
computerizată. Google MediaPipe \cite{mediapipe2019} oferă un model de estimare a
pozițiilor articulațiilor mâinii în timp real, capabil să funcționeze pe
hardware convențional. Acest lucru deschide posibilitatea construirii unor
sisteme de teleoperare accesibile, în care un operator uman controlează o mână
robotică prin simpla filmare a propriei mâini cu o cameră web.

Provocarea principală constă în integrarea celor două domenii, adică procesarea video
și extragerea coordonatelor articulațiilor, realizată pe un calculator gazdă, împreună cu
controlul determinist al servomotoarelor, realizat pe o platformă hardware dedicată. Această lucrare abordează tocmai
această integrare, propunând o soluție bazată pe FPGA (\textit{Field-Programmable
Gate Array}, matrice de porți programabile) pentru componenta de control.

\section{Motivația personală}
\label{sec:intro-motivatie-personala}

Pasiunea pentru construirea de lucruri concrete și plăcerea de a vedea un
proiect prinzând viață m-au îndreptat către o temă de licență cu un puternic
caracter aplicativ. Mi-am dorit să realizez ceva util
și, în același timp, spectaculos, iar domeniul teleoperării m-a fascinat
dintotdeauna prin ideea de a controla de la distanță un sistem mecanic care
reproduce gesturile umane. Din această dorință s-a născut tema lucrării de
față, un braț robotic care imită în timp real mișcările mâinii umane,
preluate cu ajutorul unei camere web. Componenta de viziune computerizată
este dezvoltată de colegul meu Lion Moroz, în timp ce contribuția mea
acoperă preluarea datelor prin UART (\textit{Universal Asynchronous
Receiver-Transmitter}) și prelucrarea lor pe placa Basys3, unde
sunt implementate filtrul Kalman și celelalte module de control descrise în
această lucrare. Pentru realizarea sistemului am utilizat, de asemenea, o
gamă largă de componente Analog Devices, după cum se va observa pe parcurs.

Neavând experiență în modelarea 3D, am pornit de la modele open-source
existente. Primul care mi-a atras atenția a fost brațul prezentat de Horizon
într-un material video \cite{horizon_hand}. I-am scris autorului pentru a-i
solicita fișierele, iar acesta a răspuns și le-a publicat pe platforma
Thingiverse. Deși ulterior am ales un alt model, mai potrivit cerințelor
proiectului, această interacțiune cu comunitatea open-source mi-a oferit un
impuls și o motivație suplimentară pentru a finaliza un proiect de calitate.

\section{Motivația alegerii platformei FPGA}
\label{sec:intro-fpga}

Alegerea unei platforme FPGA pentru
controlul servomotoarelor, în locul unui microcontroler tradițional, este
motivată de mai multe avantaje fundamentale.

\textbf{Determinism temporal.} Un microcontroler execută instrucțiuni secvențial,
iar întreruperile, schimbările de context și accesul la memorie introduc
variabilitate în timpul de execuție (\textit{jitter}). Într-un FPGA, logica de
control este implementată direct în hardware, iar semnalele PWM (\textit{Pulse
Width Modulation}, modulație în lățimea pulsului) sunt generate de
circuite dedicate care funcționează sincron cu ceasul de sistem. Rezultatul este
un jitter de ordinul nanosecundelor, comparativ cu microsecundele sau chiar
zecile de microsecunde ale unui microcontroler. Pentru servomotoare de tip
MG996R, care interpretează lățimea pulsului PWM ca poziție angulară, această
precizie temporală se traduce direct în stabilitatea poziționării.

\textbf{Paralelism real.} Un FPGA permite generarea simultană a tuturor celor
8 canale PWM, fără multiplexare în timp. Fiecare canal PWM dispune de propriul
comparator hardware, iar actualizarea poziției unui servo nu afectează în
niciun fel celelalte canale. În plus, modulul de filtrare Kalman, interfața
UART și controlerul SPI (\textit{Serial Peripheral Interface}) pentru ADC-ul
(\textit{Analog-to-Digital Converter}, convertor analog-digital) AD7124-8
funcționează toate în paralel,
fără a concura pentru același procesor.

\textbf{Filtrare Kalman în timp real.} Filtrul Kalman este un estimator optim
pentru sisteme liniare cu zgomot gaussian. În contextul acestei lucrări, acesta
modelează fiecare servo ca un sistem cu două stări (poziție și viteză)
și estimează traiectoria reală pe baza măsurătorilor zgomotoase primite de la
MediaPipe. Implementarea în FPGA, utilizând aritmetică în virgulă fixă Q16.16,
permite execuția completă a algoritmului într-un număr fix de cicluri de ceas,
garantând că filtrarea se finalizează înainte de următorul ciclu de servo
(20~ms).

\textbf{Scalabilitate.} Arhitectura modulară a proiectului permite adăugarea
de noi canale PWM sau senzori fără modificarea logicii existente, în limita
resurselor FPGA disponibile.

\section{Prezentarea generală a sistemului}
\label{sec:intro-overview}

Sistemul propus în această lucrare realizează controlul unei mâini robotice
printate 3D, cu 8 servomotoare MG996R, pe baza gesturilor capturate în timp
real de la un operator uman. Arhitectura de ansamblu este structurată pe
două niveluri:

\begin{enumerate}
	\item \textbf{Nivelul de achiziție}: un calculator gazdă rulează
	      biblioteca MediaPipe pentru detecția articulațiilor mâinii din fluxul
	      video. Unghiurile calculate pentru fiecare articulație sunt transmise
	      prin UART la 115200~baud, sub forma unor cadre de 10~octeți.
	\item \textbf{Nivelul de control}: placa FPGA Basys3 (Xilinx Artix-7
	      XC7A35T, 100~MHz) recepționează cadrele UART, aplică un filtru Kalman
	      cu 2~stări pentru fiecare canal, generează 8~semnale PWM independente
	      și transmite aceste semnale prin izolatori galvanici ADUM1400 către
	      servomotoare. Un ADC AD7124-8 pe 24~de biți, interfațat prin SPI
	      printr-un izolator MAX14850, permite citirea poziției reale a
	      servomotoarelor pentru verificare.
\end{enumerate}

Lanțul complet de semnal este:
\begin{center}
\mbox{\textit{Cameră} $\rightarrow$ \textit{MediaPipe} $\rightarrow$ \textit{UART} $\rightarrow$ \textit{FPGA} $\rightarrow$ \textit{ADUM1400} $\rightarrow$ \textit{Servo}}
\end{center}

Cele 8~servomotoare sunt mapate pe articulațiile mâinii robotice în ordinea
următoare, și anume degetul mic, degetul inelar, degetul mijlociu, degetul arătător,
mișcarea palmară a degetului mare, tensiunea degetului mare, rotația
încheieturii și extensia cotului. Această configurație permite reproducerea
majorității gesturilor naturale ale mâinii.

Un mod demonstrativ (\textit{demo mode}), activabil prin intermediul
comutatoarelor de pe placa Basys3, permite funcționarea autonomă a
sistemului prin parcurgerea unei secvențe predefinite de mișcări stocate
într-un modul ROM, fără a necesita conexiunea cu calculatorul gazdă.
